\chapter{Chapter06 — Reduced Vandermonde}
%%% generated by the blueprint pipeline from FABL.Chapter06.Constructions.ReducedVandermonde
\section{Overview}

Declarations of the Lean module \texttt{FABL.Chapter06.Constructions.ReducedVandermonde}.

\section{Declarations}

\begin{definition}[Reduced Vandermonde row index set]
\lean{FABL.ReducedVandermondeBinaryRow}
\statementsource{Analysis-of-Boolean-Functions-by-Ryan-ODonnell}{pdf-ffbe9c7b3306-p154-b007}
\label{FABL.ReducedVandermondeBinaryRow}
\leanok
For nonnegative integers $k$ and $\ell$, define the reduced Vandermonde binary row index
set to consist of one distinguished element together with pairs $(i,r)$ where $0\le
i<\lfloor k/2\rfloor$ and $0\le r<\ell$.
\end{definition}

\begin{definition}[Reduced Vandermonde row bijection]
\lean{FABL.reducedVandermondeBinaryRowEquiv}
\statementsource{Analysis-of-Boolean-Functions-by-Ryan-ODonnell}{pdf-ffbe9c7b3306-p154-b007}
\label{FABL.reducedVandermondeBinaryRowEquiv}
\leanok
\uses{FABL.ReducedVandermondeBinaryRow}
For nonnegative integers $k$ and $\ell$, there is a bijection between the indices $0\le
t<\lfloor k/2\rfloor\ell+1$ and the set consisting of one distinguished element together
with pairs $(i,r)$ satisfying $0\le i<\lfloor k/2\rfloor$ and $0\le r<\ell$.
\end{definition}

\begin{definition}[Binary reduced Vandermonde matrix]
\lean{FABL.reducedVandermondeBinaryMatrixOfPoints}
\statementsource{Analysis-of-Boolean-Functions-by-Ryan-ODonnell}{pdf-ffbe9c7b3306-p154-b007}
\label{FABL.reducedVandermondeBinaryMatrixOfPoints}
\leanok
\uses{FABL.BinaryExtensionField, FABL.binaryExtensionEncode, FABL.reducedVandermondeBinaryRowEquiv, FABL.𝔽₂}
For natural numbers $k,n$ and a positive extension degree $\ell$, let
$\alpha_1,\ldots,\alpha_n$ be elements of the binary extension field of degree $\ell$
over $\mathbb{F}_2$. Define a matrix over $\mathbb{F}_2$ with $n$ columns and $(\lfloor
k/2\rfloor\ell+1)$ rows: its first row is constant $1$, and its remaining rows are the
binary coordinates of the odd powers $\alpha_j^{2q+1}$, for each applicable exponent
index $q$ and each coordinate in the extension field representation.
\end{definition}

\begin{theorem}[Constant row of the binary Vandermonde matrix]
\lean{FABL.reducedVandermondeBinaryMatrixOfPoints_constantRow}
\proofsource{Analysis-of-Boolean-Functions-by-Ryan-ODonnell}{pdf-ffbe9c7b3306-p154-b007}
\label{FABL.reducedVandermondeBinaryMatrixOfPoints_constantRow}
\leanok
\uses{FABL.BinaryExtensionField, FABL.reducedVandermondeBinaryMatrixOfPoints, FABL.reducedVandermondeBinaryRowEquiv}
\difficulty{1}
Let $k$ and $n$ be nonnegative integers, let $ell$ be a nonzero natural number, and let
$alpha$ assign to each of the $n$ column indices an element of the degree-$\ell$
extension field of $\mathbb{F}_2$. For every column $j$, the matrix row corresponding to
the distinguished constant-row index has entry $1$.
\end{theorem}

\begin{theorem}[Odd-power rows of the binary Vandermonde matrix]
\lean{FABL.reducedVandermondeBinaryMatrixOfPoints_oddPowerRow}
\label{FABL.reducedVandermondeBinaryMatrixOfPoints_oddPowerRow}
\leanok
\uses{FABL.BinaryExtensionField, FABL.binaryExtensionEncode, FABL.reducedVandermondeBinaryMatrixOfPoints, FABL.reducedVandermondeBinaryRowEquiv}
\difficulty{1}
Let $k,n$ be nonnegative integers, let $\ell$ be positive, and let $\alpha$ assign a
point of the degree-$\ell$ extension field of $\mathbb{F}_2$ to each of $n$ columns. For
every row indexed by $q<k/2$ and a basis coordinate $i<\ell$, the entry in column $j$ is
the $i$th coordinate, in the chosen basis, of $\alpha_j^{2q+1}$.
\end{theorem}

\begin{theorem}[Vanishing of all low order moments]
\lean{FABL.extensionField_moments_eq_zero_of_oddMoments_eq_zero}
\label{FABL.extensionField_moments_eq_zero_of_oddMoments_eq_zero}
\leanok
\uses{FABL.BinaryExtensionField}
\difficulty{5}
Let $\ell,n,k$ be nonnegative integers, let $\alpha_0,\ldots,\alpha_{n-1}$ lie in the
finite field with $2^\ell$ elements, and let $S\subseteq\{0,ldots,n-1\}$. Suppose that
$|S|=0$ in this field and that, for every nonnegative integer $q$ with $2q+1<k$, the sum
$\sum_{j\in S}\alpha_j^{2q+1}$ is zero. Then for every nonnegative integer $p<k$, the
sum $\sum_{j\in S}\alpha_j^p$ is zero.
\end{theorem}

\begin{theorem}[Reduced vandermonde binary matrix of points has nonzero column sums up to]
\lean{FABL.reducedVandermondeBinaryMatrixOfPoints_hasNonzeroColumnSumsUpTo}
\statementsource{Analysis-of-Boolean-Functions-by-Ryan-ODonnell}{pdf-ffbe9c7b3306-p154-b007}
\label{FABL.reducedVandermondeBinaryMatrixOfPoints_hasNonzeroColumnSumsUpTo}
\leanok
\uses{FABL.BinaryExtensionField, FABL.HasNonzeroColumnSumsUpTo, FABL.binaryExtensionEncode, FABL.extensionField_moments_eq_zero_of_oddMoments_eq_zero, FABL.matrixColumnSum_apply_eq_sum, FABL.reducedVandermondeBinaryMatrixOfPoints, FABL.reducedVandermondeBinaryMatrixOfPoints_oddPowerRow, FABL.reducedVandermondeBinaryRowEquiv, FABL.vandermondeBinaryMatrixOfPoints, FABL.vandermondeBinaryMatrixOfPoints_hasNonzeroColumnSumsUpTo, FABL.vandermondeBinaryRowEquiv, FABL.𝔽₂}
\difficulty{5}
% dataset-pipeline: stub -- prose pending build_informal_statements.py
Statement pending regeneration.
\end{theorem}

\begin{theorem}[Exists k wise independent subspace card le two mul two n pow half]
\lean{FABL.exists_kWiseIndependentSubspace_card_le_two_mul_two_n_pow_half}
\statementsource{Analysis-of-Boolean-Functions-by-Ryan-ODonnell}{
  pdf-ffbe9c7b3306-p154-b006,
  pdf-ffbe9c7b3306-p154-b007,
  pdf-ffbe9c7b3306-p155-b001
}
\label{FABL.exists_kWiseIndependentSubspace_card_le_two_mul_two_n_pow_half}
\leanok
\uses{FABL.BinaryExtensionField, FABL.HasNonzeroColumnSumsUpTo, FABL.IsLowDegreeFourierRegular, FABL.binaryExtensionEnumeration, FABL.binaryFunctionOnSignCube, FABL.card_submodule_eq_two_pow_finrank, FABL.matrixRowSpan, FABL.matrixRowSpanDensity, FABL.matrixRowSpanDensity_isKWiseIndependent_iff, FABL.reducedVandermondeBinaryMatrixOfPoints, FABL.reducedVandermondeBinaryMatrixOfPoints_hasNonzeroColumnSumsUpTo, FABL.subsetDensity, FABL.𝔽₂}
\difficulty{5}
% dataset-pipeline: stub -- prose pending build_informal_statements.py
Statement pending regeneration.
\end{theorem}
